\section{Introduzione}

L'informazione genetica costituisce il principale meccanismo attraverso cui le caratteristiche ereditarie vengono conservate e trasmesse
tra generazioni. Con la scoperta della struttura del DNA e, successivamente, con lo sviluppo delle tecniche di sequenziamento, è diventato
possibile osservare tale informazione non soltanto come fenomeno biologico, ma anche come una sequenza discreta di simboli.

Dal punto di vista computazionale, infatti, un filamento di DNA può essere rappresentato come una stringa definita su un alfabeto finito
in cui ciascun simbolo rappresenta una delle quattro basi azotate. Questa rappresentazione permette di trattare l'informazione genetica
come informazione elaborabile computazionalmente, trasformando problemi di natura biologica in problemi algoritmici su stringhe.

Un problema fondamentale consiste nel confrontare due sequenze biologiche allo scopo di individuarne le corrispondenze e le differenze.
Sequenze omologhe possono infatti differire a causa di mutazioni. In ambito computazionale tali differenze possono
essere studiate mediante un allineamento di sequenze, che mette in corrispondenza i simboli delle due stringhe introducendo, quando necessario,
opportuni simboli di gap. Associando a ciascuna possibile corrispondenza un valore attraverso un modello valutativo, il confronto può essere formulato
come un problema di ottimizzazione: il Global Alignment Problem, che consiste nel determinare un allineamento avente score massimo.

Una soluzione classica a questo problema è fornita dall'algoritmo di Needleman--Wunsch, basato sulla programmazione dinamica.
Tale approccio consente di ricostruire un allineamento ottimo, ma nella sua formulazione standard richiede uno spazio proporzionale al prodotto
delle lunghezze delle due sequenze.
Per input di grandi dimensioni, il consumo di memoria può quindi diventare un limite pratico significativo.

L'obiettivo di questa tesi è analizzare il problema dell'allineamento globale con particolare attenzione all'utilizzo della memoria, confrontando
l'algoritmo di Needleman--Wunsch con l'algoritmo di Hirschberg, che attraverso una strategia divide et impera permette di ricostuire un allineamento
ottimo utilizzando spazio lineare. I due algoritmi vengono analizzati teoricamente, implementati in linguaggio C e infine confrontati sperimentalmente
su sequenze sintetiche e biologiche reali, valutandone tempo di esecuzione e consumo di memoria.
